% Include fragment: requires amsmath, amssymb, amsthm and a lemma environment.
\section{Exact source and tape binding for MP196}

This section concerns the delivered 196-state prototype. It does not
assert a reduction below 194 states. The baseline is the primary PA194
table at repository commit
\texttt{2a88cb6a3f823e0de5addf19d725bf8373860c71}.

Write $p(a,b)=(a+b)(a+b+1)/2+b$ and $I(a,b)=4p(a,b)$.
Logical registers $1,2,4,5,6,7,8$ hold the accumulator $T$, formula stack
$W$, offset $D$, cache $H$, scratches $X,Y$, and command residual $L$.
The logical-to-physical permutation remains
\[
 (6,2,0,4,3,7,8,5,1).
\]
All registers not mentioned in a block's postcondition are framed.

\subsection{The unchanged operand and the new successful exit}

The original MP operand block is retained. For input command store
$(t,w,d,h)$ and parameter $x$, put $b=d+x$. The inherited operand block
reads $x$, sets $D=b$, constructs $I(t,b)$, and restores the old
formula-stack tail. The following original clear-and-compare block
compares $I(t,b)$ with $h$. It leaves $T=H=X=Y=0$ and returns directly
to the command loop on inequality. Thus failure returns $(0,w,b,0)$.
On equality it enters source site 3634 with $D=b$ and the same cleared
registers.

Only the successful source continuation is replaced. In the structured
source, the precise replacement is
\[
 \mathtt{node163}
 =\bigl(\mathtt{node6},
     (\mathtt{node51},(\mathtt{inc}\,5,\mathtt{node7}))\bigr).
\]
The original definition
$\mathtt{node7}=(\mathtt{node6},\mathtt{halt})$ is retained, including
its original occurrence in Shift. The transfer block \texttt{node51}
is also retained; it has the exact numerical instructions
\[
\begin{array}{rll}
3634&\mathtt{dec}\,D&0\mapsto3638,\quad +\mapsto3635,\\
3635&\mathtt{inc}\,T&\mapsto3636,\\
3636&\mathtt{inc}\,X&\mapsto3637,\\
3637&\mathtt{jump}&\mapsto3634,\\
3638&\mathtt{inc}\,H&\mapsto3639.
\end{array}
\]
A positive decrement subtracts one. At the head of the transfer loop,
the invariants $T+D=b$ and $X+D=b$ hold, while $H=Y=0$.
Every positive iteration decreases $D$. It therefore reaches site 3638
with $T=X=b$ and $D=H=Y=0$. The increment sets $H=1$.

The inherited destructive comparator is now entered at site 3639 with
$H=1$ and $X=b$. Its complete new occurrence has the CFG
\[
\begin{array}{rll}
3639&\mathtt{dec}\,H&0\mapsto3643,\quad +\mapsto3640,\\
3640&\mathtt{dec}\,X&0\mapsto3642,\quad +\mapsto3641,\\
3641&\mathtt{jump}&\mapsto3639,\\
3642&\mathtt{jump}&\mapsto3644,\\
3643&\mathtt{dec}\,X&0\mapsto3649,\quad +\mapsto3644,\\
3644&\mathtt{dec}\,H&0\mapsto3646,\quad +\mapsto3645,\\
3645&\mathtt{jump}&\mapsto3644,\\
3646&\mathtt{dec}\,X&0\mapsto3648,\quad +\mapsto3647,\\
3647&\mathtt{jump}&\mapsto3644,\\
3648&\mathtt{jump}&\mapsto3650,\\
3649&\mathtt{halt},&\\
3650&\mathtt{jump}&\mapsto25.
\end{array}
\]
If $b=0$, site 3640 takes its zero branch and the cleanup returns to
site 25 with $T=0$. If $b=1$, the first pair of positive decrements
sets $H=X=0$, after which sites 3639 and 3643 take their zero branches
to halt. If $b\ge2$, those branches leave $X=b-2$ at site 3644; the
cleanup decreases $X$ to zero and returns to site 25. In every
nonhalting case $T=b$, $D=H=X=Y=0$, and the formula stack is unchanged.
Consequently the exact MP command is
\[
 \operatorname{MP}_x(t,w,d,h)=
 \begin{cases}
  (0,w,d+x,0),&h\ne I(t,d+x),\\
  \mathsf{halt},&h=I(t,d+x)\ \text{and}\ d+x=1,\\
  (d+x,w,0,0),&h=I(t,d+x)\ \text{and}\ d+x\ne1.
 \end{cases}
\]
The finite builder asserts all the displayed control endpoints, so the
new successful continuation cannot accidentally return to Shift instead
of the command loop. The MP source entry is 3554 and the Shift source
entry is 3651.

\subsection{The retained Shift and a syntactic invariant}

The original Shift and its halt leaf are unchanged. On an arbitrary
command store, Shift halts when $t+h=1$, and otherwise produces
$(h,\operatorname{tl}(w),d,\operatorname{hd}(w))$.
The old halt leaf remains source site 10, path $001$, and is still
covered by its primitive tape contract.

Define the canonical contents of a numerical stack by repeatedly
applying the Cantor projections until its tail is zero, then regarding
further heads as zero. This is finite because for $w>0$,
$\operatorname{tl}(w)<w$, while both projections of zero are zero.
Say a command store is \emph{one-free} when $t$, $h$, and every canonical
stack entry differ from 1.

\begin{lemma}
Starting at the initialized command store, every completed nonhalting
command boundary is one-free, without assuming consistency of PA.
\end{lemma}
\begin{proof}
Initially the formula fields and formula stack are zero. The unchanged
AX command pushes the saved top and nine generated formulas: the
generated implication heads and gate are divisible by four, and the
universal head is congruent to two and at least two. None equals one.
AX resets the cache to zero. An unsuccessful MP returns top and cache
zero. A successful MP with consequent one halts without returning to a
completed command boundary; every other successful MP stores zero or
a value at least two. MP
leaves the formula stack unchanged. Finally, if nonnegative $t,h$
are both different from one, then $t+h\ne1$. Shift therefore takes
its ordinary nonhalting branch and merely exposes existing one-free
stack contents. Induction over the finite sequence of completed commands proves
the claim.
\end{proof}

The successful MP transfer may transiently write $T=1$ inside a command;
the invariant is asserted only at completed nonhalting command boundaries.
Thus the original Shift halt is unreachable after initialization. The
abstract protocol with total Shift and the displayed MP rule simulates
every literal execution from that frame up to and including its first
halt. The command proof compiler uses MP parameter $b$ for conclusion
$b$, and may discard the cleared successful offset immediately with
Pop. This supplies the connection to the unshifted immediate-MP-halt
protocol in the main report. No consistency assumption is used to
exclude the old Shift halt.

\subsection{Literal controls and exact startup}

The only changed inherited row is $(164,1)$, which becomes $1R400$.
The two new readers have the entire transition table
\[
\begin{array}{c|cc}
 &0&1\\\hline
400&0R130&1R401\\
401&0R63&1R68.
\end{array}
\]
Native controls 130, 63, and 68 implement the existing transfer block,
increment of $H$, and comparator-plus-halt block in these source
contexts. The 23 controls used by the inherited arithmetic-block
oracle have exactly their original rows.
The dense machine has 196 working controls, a separate halt label 196,
and 392 binary left/right rows. Every working control is reachable in
the transition graph from the initial control.

The verifier resolves 3,663 source sites and checks all 4,768 primitive
outcomes against symbolic frames, including both outcomes of every
decrement and both occurrences of the halt instruction. It uses
1,476,043 symbolic operations, with maximum 497 operations for one
primitive contract. The simulation induction behind these checks
inherits the unchanged generic arithmetic-block lemmas and uses
explicit one-row steps outside them. These executable certificates
are not an independent proof-assistant kernel theorem or a new
literal verification of those arithmetic blocks.

A separate literal blank-tape simulation takes exactly 58,154
transitions and checks equality of the entire finite tape support
with a clean frame: native state 84, head $-2$, source site 25, gap 48,
nine physical primary registers
\[
 (0,0,0,0,0,0,2,0,0),
\]
and 26 further registers of value one. The whole tape has 66 one-cells.
This represents candidate $K=2$ and all other logical registers zero.
It lies in the symbolic frame scope with gap parameter $n=0$.
Every initialization transition is checked to avoid the halt label.

The ordinary PA-formalizable command proof supplied in the main report
depends on the attributed numerical-calculus calibration from the
existing packets. There is no new end-to-end Lean/Rocq theorem or
claimed iPRA internalization. The construction remains two states
larger than the primary 194-state baseline.
